% !TEX TS-program = lualatex
% !TEX encoding = UTF-8

\documentclass[VesperaleOP.tex]{subfiles}

\ifcsname preamble@file\endcsname
  \setcounter{page}{\getpagerefnumber{M-vop53_sanctorale_martii}}
\fi

\begin{document}

\festum{}{FESTA MARTII}{}{}{2}
\festum{4.}{S. Casimiri Confessoris}{III classis}{festa martii. 4.}{3}
\oratio{Deus, qui inter regáles delícias et mundi illécebras sanctum Casimírum virtúte constántiæ roborásti:\psstar{} quǽsumus ut eius intercessióne fidéles tui terréna despíciant,\pscross{}et ad cæléstia semper aspírent.\psstar{} Per.}
\lineamagna

%%%%
\festum{die 7 martii}{S. Thomæ de Aquino\\Confessoris O. P. et Ecclesiæ Doctoris}{I classis}{festa martii. 7.}{1}
\titulumrubrica{Ad I Vesperas}
\cantus[Ant. I b]{0306VA}{Super Ps.}
\capitulum{Eccli. 45}{Elégit eum Dóminus ex omni carne,\pscross{}et dedit illi coram præcépta et legem vitæ et disciplínæ,\psstar{} docére Iacob testaméntum suum et iudícia sua Israël.}
\cantus{0306VR}{\rrann}
\cantus{0306VH}{Hymnus}
\cantus{0306VAM}{Ad Magn.}
\oratio{Deus, qui Ecclésiam tuam beáti Thomæ Confessóris tui atque Doctóris mira eruditióne claríficas et sancta operatióne fecúndas:\psstar{} da nobis, quǽsumus, et quæ dócuit intelléctu conspícere,\pscross{}et quæ egit imitatióne complére.\psstar{} Per.}

\rubrica{Si venerit in Quadragesima, fit commemoratio feriæ.}

\titulumrubrica{Ad II Vesperas}
\cantus[Ant. I a]{0307VA}{Super Ps.}
\rubrica{Psalmi de Communi unius Confessoris \normaltext{???}}
\rubrica{Capitulum, hymnus et \vvrub ut in I Vesperis \normaltext{\pageref{M-0306VR}.}}
\cantus{0307VAM}{Ad Magn.}
\rubrica{Oratio ut in I Vesperis et fit commemoratio feriæ tantum, si venerit in Quadragesima.}
\lineamagna

%%%%
\festum{8.}{S. Ioannis a Deo Confessoris}{III classis}{festa martii. 8.}{3}
\oratio{Deus, qui beátum Ioánnem, tuo amóre succénsum, inter flammas innóxium incédere fecísti,\pscross{}et per eum Ecclésiam tuam nova prole fecundásti:\psstar{} præsta, ipsíus suffragántibus méritis, ut igne caritátis tuæ vítia nostra curéntur,\pscross{}et remédia nobis ætérna provéniant.\psstar{} Per.}
\lineamagna

%%%%
\festum{9.}{S. Franciscæ Romanæ}{III classis}{festa martii. 9.}{3}
\oratio{Deus, qui beátam Francíscam fámulam tuam, inter cétera grátiæ tuæ dona, familiári Angeli consuetúdine decorásti:\psstar{} concéde, quǽsumus, ut intercessiónis eius auxílio\pscross{}Angelórum consórtium cónsequi mereámur.\psstar{} Per.}
\lineamagna

%%%%
\festum{die 19 martii}{S. Ioseph, Sponsi B. Mariæ V.}{Confessoris et Ecclesiæ universalis Patroni\\I classis}{festa martii. 19.}{1}
\titulumrubrica{Ad I Vesperas}
\cantus[Ant. VIII a]{0318VA}{Super Ps.}
\capitulum{Prov. 28 et 27}{Vir fidélis multum laudábitur:\psstar{} et qui custos est dómini sui glorificábitur.\pscross{}}
\cantus{0318VR}{\rrann}
\cantus[1. cantus II]{0318VH1}{Hymnus}
\cantus[2. cantus I]{0318VH2}{Hymnus}
\cantus[3. cantus VII]{0318VH3}{Hymnus}

\versiculus{Amávit eum Dóminus, et ornávit eum. Allelúia.}{Stolam glóriæ índuit eum. Allelúia.}

\cantus{0318VAM}{Ad Magn.}
\oratio{Beáti Patriárchæ Ioseph, sanctíssimæ Genetrícis tuæ Sponsi, quǽsumus, Dómine, méritis adiuvémur:\psstar{} ut, quod possibílitas nostra non óbtinet,\pscross{}eius nobis intercessióne donétur:\psstar{} Qui vivis.}

\rubrica{Si venerit in Quadragesima, fit commemoratio feriæ.}

\titulumrubrica{Ad II Vesperas}
\cantus[Ant. VI]{0319VA}{Super Ps.}
\rubrica{Psalmi de Communi unius Confessoris \normaltext{???}}

\rubrica{Capitulum, hymnus et \vvrub ut in I Vesperis \normaltext{\pageref{M-0318VH1}.}}

\cantus{0319VAM}{Ad Magn.}
\rubrica{Oratio ut in I Vesperis et fit commemoratio feriæ tantum, si venerit in Quadragesima.}
\lineamagna

%%%%
\festum{die 25 martii}{In Annuntiatione B. M. V.}{I classis}{festa martii. 25.}{1}
\titulumrubrica{Ad I Vesperas}
\cantus[Ant. I b]{0324VA}{Super Ps.}
\capitulum{Is. 7}{Ecce Virgo concípiet et páriet fílium,\pscross{}et vocábitur nomen eius Emmánuel.\psstar{} Butýrum et mel cómedet, ut sciat reprobáre malum et elígere bonum.\pscross{}}
\cantus{0324VR}{\rrann}
\cantus{0324VH}{Hymnus}

\versiculus{Roráte, cæli, désuper, et nubes pluant iustum. Allelúia.}{Aperiátur terra, et gérminet Salvatórem. Allelúia.}

\cantus{0324VAM}{Ad Magn.}
\oratio{Deus, qui de beátæ Maríæ Vírginis útero Verbum tuum, Angelo nuntiánte, carnem suscípere voluísti,\pscross{}præsta supplícibus tuis:\psstar{} ut, qui vere eam Genetrícem Dei crédimus,\pscross{}eius apud te intercessiónibus adiuvémur.\psstar{} Per eúndem.}

\rubrica{Si venerit in Quadragesima, fit commemoratio feriæ.}

\titulumrubrica{Ad II Vesperas}
\cantus[Ant. VIII a]{0325VA}{Super Ps.}
\rubrica{Psalmi ut in festis B. Mariæ Virg. \normaltext{???}}
\rubrica{Capitulum, hymnus et \vvrub ut in I Vesperis \normaltext{\pageref{M-0324VA}.}}
\cantus{0325VAM}{Ad Magn.}
\rubrica{Oratio ut in I Vesperis et fit commemoratio feriæ tantum, si venerit in Quadragesima.}
\lineamagna

%%%%
\festum{feria sexta post dominicam in passione}{In Compassione B. Mariæ Virg.}{II classis}{In compassione b. m. v.}{2}
\cantus[Ant. VII b]{0331VA}{Super Ps.}
\capitulum{Is. 21}{Angústia possédit me sicut angústia parturiéntis:\psstar{} conturbáta sum cum vidérem:\pscross{}emárcuit cor meum, ténebræ stupefecérunt me.}
\cantus{0331VH}{Hymnus}

\versiculus{Sancta Dei Génitrix, Virgo dulcis atque pia.}{Christum Regem morti tráditum pro nobis exóra.}

\cantus{0331VAM}{Ad Magn.}
\oratio{Intervéniat pro nobis, quǽsumus, Dómine Iesu Christe, nunc et in hora mortis nostræ, apud tuam cleméntiam beáta Virgo María Mater tua:\psstar{} cuius sacratíssimam ánimam in hora benedíctæ Passiónis tuæ dolóris gládius pertransívit,\pscross{}et in gloriósa Resurrectióne tua ingens gáudium lætificávit:\psstar{} Qui vivis.}

\rubrica{Commemoratio feriæ.}
\lineamagna

\end{document}


